\subsection{Amount of tissue}
\label{sec:amount-methods}
For each suction we sum three measures over its \texttt{dyed} and \texttt{solid\_no\_dye}
pieces: the number of pieces, height\,$\times$\,width, and the area (counts\,$\times$\,ms).
We chose the area as the main measure before seeing the results, because it is the fluorescence
added up over time.

Let $S_{ij}$ be one of these measures for suction $j$ (first suction or re-suction) of
recording $i$, and $A_i$ the known amount aspirated in that recording (ml or g). For each group of
recordings we fitted the linear mixed model
\begin{equation}
  \log S_{ij} = \log k + \log A_i + \beta_j + u_i + \varepsilon_{ij},
  \label{eq:amount}
\end{equation}
where $\beta_j$ is the effect of the suction number and $u_i$ a random effect of the recording.
% (R: \texttt{lmer(log(S) $\sim$ offset(log(A)) + pass + (1 | rec))})
The amount of a first
suction is estimated as $\hat{A} = S / k$.

We also tried measures that do not need the concentration. The height and the area both grow with the dye
concentration, so their ratio should not. We took the sum over the pieces of area\,/\,height,
similar to pulse-width sizing in flow cytometry~\cite{sharpless1976}, and of two widths measured
relative to each piece's own height (full width at half maximum of the highest peak, and the
time above half of the largest peak).

The application does not show the amount, because the area per millilitre depends on the
concentration.
